% ===================================================================
\section{Non-Geometric Topologies}
\label{sec:nongeometric}

Section~\ref{sec:programmability} established that the Steersman
requires geometric coherence for block-diagonal structure to emerge.
This section examines the two non-geometric experiments in detail:
circular resonance (R-001, number-theoretic pairs) and the emanation
architecture (E-001, hierarchical bridge structure). Together with
WL-001 (wrong-labels), they define the boundary between successful
and unsuccessful topology programming.

\subsection{Circular Resonance: The Strongest Negative (R-001)}

The resonance topology derives co-axial pair assignments from
number-theoretic relationships: Sophie Germain pairing
($2(11) + 1 = 23 \to$ channels $(0, 2)$), consecutive-prime pairing
($29, 31 \to$ channels $(3, 4)$), and residue pairing
($11 \equiv 17 \equiv 5 \pmod{6} \to$ channels $(1, 5)$). Unlike
the RD pairs, which correspond to opposite faces of a Voronoi cell,
the resonance pairs have algebraic but not spatial motivation.

\begin{table}[H]
\centering
\caption{Resonance control (R-001) vs.\ wrong-labels (WL-001) vs.\
  RD contrastive (H-ch6). All three use $n = 6$ with identical
  hyperparameters; only the pair specification differs.}
\label{tab:resonance}
\small
\begin{tabular}{lcccccc}
\toprule
Experiment & Pair Source & Steps & $\rho$ & Fiedler & Deviation & BD \\
\midrule
H-ch6 (RD) & Spatial & 10{,}000 & 70{,}404:1 & $9 \times 10^{-5}$ & --- & Yes \\
WL-001 (random) & Random & 10{,}000 & $8.7 \times 10^{-6}$ & $1.3 \times 10^{-5}$ & 2.12 & No \\
R-001 (primes) & Algebraic & 10{,}000 & $8.9 \times 10^{-6}$ & $1.2 \times 10^{-5}$ & 2.12 & No \\
\bottomrule
\end{tabular}
\end{table}

R-001 is the strongest negative control in the programme. The
co-axial/cross-axial coupling ratio decayed monotonically from
$0.257$ at step~100 to $6.5 \times 10^{-4}$ at step~1{,}000 and
$8.9 \times 10^{-6}$ at step~10{,}000. Fiedler follows the same
trajectory: $1.7 \times 10^{-4}$ at step~1{,}000, $1.2 \times 10^{-5}$
at step~10{,}000. Deviation from identity grows throughout training
to~2.12, indicating that the bridge matrix departs substantially
from identity without developing any directional structure---an
extreme, unorganized coupling pattern.

Validation loss converges normally ($0.4008$ at step~10{,}000),
confirming that the resonance topology does not impede language
modelling. The adapter learns the task while the bridge fails to
organize.

\subsection{WL-001 and R-001: Functional Equivalence}

A key finding emerges from comparing the two negative controls:
WL-001 (random partition) and R-001 (number-theoretic pairs) produce
\emph{nearly identical trajectories} (Figure~\ref{fig:wl_r001_identity}).
Table~\ref{tab:wl_r_compare} shows matched checkpoints.

\begin{table}[H]
\centering
\caption{Trajectory comparison: WL-001 (random partition) vs.\
  R-001 (number-theoretic pairs). Maximum divergence ${<}\,10\%$
  at all checkpoints.}
\label{tab:wl_r_compare}
\small
\begin{tabular}{ccccc}
\toprule
Step & WL-001 $\rho$ & R-001 $\rho$ & WL-001 Fiedler & R-001 Fiedler \\
\midrule
1{,}000 & $6.5 \times 10^{-4}$ & $7.1 \times 10^{-4}$ & $1.6 \times 10^{-4}$ & $1.7 \times 10^{-4}$ \\
5{,}000 & $5.0 \times 10^{-5}$ & $5.1 \times 10^{-5}$ & $6.0 \times 10^{-5}$ & $6.4 \times 10^{-5}$ \\
10{,}000 & $8.7 \times 10^{-6}$ & $8.9 \times 10^{-6}$ & $1.3 \times 10^{-5}$ & $1.2 \times 10^{-5}$ \\
\bottomrule
\end{tabular}
\end{table}

The maximum divergence between WL-001 and R-001 is less than 10\%
at all checkpoints. From the Steersman's perspective,
number-theoretic structure without spatial embedding is
\emph{functionally equivalent to random noise}. The algebraic
relationships (Sophie Germain, twin primes, residue classes) carry
no information that the contrastive loss can use to produce
directional coupling.

This equivalence has a precise meaning. The contrastive loss
operates on the absolute magnitudes of bridge entries at the
indices specified by the pair partition. For the contrastive gradient
to produce block-diagonal structure, the specified co-axial entries
must be amenable to selective strengthening relative to the
cross-axial entries. When the pair specification reflects the
geometry of a polytope (RD, tesseract, octahedron), the bridge's
gradient landscape has a clear minimum corresponding to the
block-diagonal attractor. When the pair specification is algebraic
(resonance) or random (wrong-labels), no such minimum exists---the
gradient landscape is flat with respect to the co-axial/cross-axial
distinction, and the contrastive pressure distributes uniformly,
driving total collapse.

Both experiments converge to Regime~4 (connectivity collapse):
Fiedler $\approx 10^{-5}$, co/cross $\approx 10^{-5}$, deviation
$\approx 2.1$. The bridge develops extreme but undirected coupling
and then decouples. This is distinct from the spectral attractor
(Regime~3), where the bridge maintains stable connectivity at
Fiedler~$\approx 0.09$. The contrastive loss, given incoherent
pairs, overwhelms the spectral loss's connectivity target and
suppresses all coupling---a destructive interaction between the
two loss components.

\begin{figure}[t]
  \centering
  \includegraphics[width=\textwidth]{figures/fig_wl_r001_identity.pdf}
  \caption{Functional equivalence of wrong-labels (WL-001, random
    partition) and resonance (R-001, number-theoretic pairs). Both
    collapse trajectories are indistinguishable ($<10\%$ divergence
    across 10{,}000 steps), demonstrating that algebraic structure
    without spatial embedding carries no information the contrastive
    loss can use.}
  \label{fig:wl_r001_identity}
\end{figure}

\subsection{Emanation Architecture: The Third Regime (E-001)}
\label{sec:emanation_results}

The emanation architecture introduces a qualitatively different
outcome. Using a master bridge matrix with per-layer offsets scaled
by $1/\sqrt{l}$ (Eq.~\ref{eq:emanation}), the same RD-derived
co-axial pairs are specified, but the contrastive loss operates
through the master matrix rather than per-layer bridges directly.

\begin{table}[H]
\centering
\caption{Emanation results (E-001) at 10{,}000 steps vs.\ standard
  RD contrastive (H-ch6) and spectral attractor (H-ch3).}
\label{tab:emanation}
\small
\begin{tabular}{lcccccl}
\toprule
Experiment & Architecture & $\rho$ & Fiedler & Deviation & Val Loss & Regime \\
\midrule
H-ch6 (RD) & Independent & 70{,}404:1 & $9 \times 10^{-5}$ & --- & 0.4015 & BD \\
E-001 (emanation) & Hierarchical & 1.12:1 & 0.084 & 0.35 & 0.4009 & Hierarchical \\
H-ch3 (spectral) & Independent & N/A & 0.095 & --- & 0.4020 & Attractor \\
\bottomrule
\end{tabular}
\end{table}

E-001 occupies a regime distinct from both block-diagonal emergence
and the spectral attractor. Its final metrics at 10{,}000 steps:

\begin{itemize}
  \item \textbf{Fiedler: 0.084.} Near but below the spectral attractor
    ($0.09$--$0.10$). The bridge maintains substantial connectivity.
  \item \textbf{Co/cross: 1.12:1.} Slight directional preference---the
    bridge favors the RD co-axial pairs by 12\% on average. This is
    above unity but far below the 70{,}000:1 achieved by the standard
    architecture with the same pair specification.
  \item \textbf{Deviation: 0.35.} The per-layer offsets have grown
    substantially from their initialized values, indicating that
    individual layers differentiate from the master bridge. This
    deviation is roughly one-sixth the deviation observed in the
    collapse regime (2.12 for WL-001/R-001).
  \item \textbf{Val loss: 0.4009.} Marginally better than both the
    standard contrastive (0.4015) and spectral-only (0.4020)
    experiments, though the difference is within noise.
\end{itemize}

\paragraph{Interpretation.}
The emanation architecture creates a tension between global coherence
and local adaptation. The master bridge receives the contrastive
gradient directly, developing slight co-axial preference. But the
per-layer offsets dilute this preference: each layer adapts its bridge
to local gradient signals, and these local adaptations do not
consistently reinforce the global topology. The master bridge carries
a weak geometric signal (12\% co-axial preference); the offsets
fragment it.

The result is a bridge that is \emph{coherent but not structured}:
it maintains connectivity (Fiedler near the attractor), develops weak
directional preference (co/cross slightly above unity), and allows
per-layer differentiation (growing deviation)---without achieving the
dramatic block-diagonal separation that requires per-bridge
contrastive specificity.

This hierarchical coherence regime suggests that topology
programming requires direct contrastive feedback to individual
bridge matrices. Indirection through a shared master bridge preserves
connectivity but loses the topological specificity needed for
block-diagonal emergence. The emanation architecture may have value
for other purposes---regularization, transfer, compositional
adapters---but it does not serve as a topology programmer.

\subsection{Topology Taxonomy: Four Regimes}

The complete experimental programme produces four distinct dynamical
regimes, cleanly separated with no intermediate or overlapping
outcomes.

\begin{enumerate}
  \item \textbf{Block-diagonal} (geometric contrastive).
    Polytope-derived pair specifications with per-bridge contrastive
    loss. Co/cross: $42{,}000$--$474{,}000$:1. Fiedler: $10^{-5}$--$10^{-4}$.
    Experiments: O-001, H-ch6, Seed-43, T-001r2.

  \item \textbf{Spectral attractor} (no contrastive).
    Spectral regularization alone, no pair specification.
    Co/cross: ${\sim}1$:1. Fiedler: $0.092$--$0.102$.
    Experiments: H-ch3, H-ch4, H-ch8, H-ch12.

  \item \textbf{Hierarchical coherence} (emanation architecture).
    Contrastive loss on shared master bridge with per-layer offsets.
    Co/cross: $1.12$:1. Fiedler: $0.084$.
    Experiment: E-001.

  \item \textbf{Connectivity collapse} (incoherent contrastive).
    Non-spatial pair specifications (random or algebraic).
    Co/cross: ${\sim}10^{-5}$:1. Fiedler: ${\sim}10^{-5}$.
    Deviation: ${\sim}2.1$.
    Experiments: WL-001, R-001.
\end{enumerate}

Two features of this taxonomy are notable. First, the collapse regime
(Regime~4) is distinct from the spectral attractor (Regime~3). The
spectral attractor is a near-initialization resting state at Fiedler
$\approx 0.09$---inside the BM-000 identity-plus-noise null, not a
distinct structured equilibrium; the collapse regime represents a
destructive interaction where the contrastive loss overwhelms the
spectral loss and drives connectivity to zero. The spectral attractor
is the bridge's resting state; the collapse regime is the bridge's
failure mode under incoherent programming.

Second, Regimes~2 and~3 are structurally similar (both near
Fiedler~$\approx 0.09$, both with co/cross near unity) but
mechanistically different. The spectral attractor sits at
near-initialization connectivity under spectral loss alone.
Hierarchical coherence reaches a similar Fiedler through a different
pathway---master bridge with contrastive loss---and develops weak
directionality that the spectral attractor does not. The regimes are adjacent in
metric space but distinct in mechanism.
